% Report skeleton for DD2430 Group 3: Agentic AI Compiler (Ericsson).
% Target: 6-10 pages, IEEE two-column conference format.
% Each section lists what belongs there and a rough page budget.
% Delete the guidance comments and \todo notes as the text gets written.

\documentclass[conference]{IEEEtran}

\usepackage[utf8]{inputenc}
\usepackage[T1]{fontenc}
\usepackage{amsmath}
\usepackage{graphicx}
\usepackage{booktabs}
\usepackage{xcolor}
\usepackage{listings}
\usepackage{hyperref}

\newcommand{\todo}[1]{\textcolor{red}{[TODO: #1]}}

\lstset{basicstyle=\footnotesize\ttfamily, breaklines=true, frame=single}

\begin{document}

\title{An Agentic Optimization Loop Around the XLA Compiler\\
\todo{final title}}

\author{
\IEEEauthorblockN{Joel Maharena, John Christensen, Lidya Nasser, Sara Strandberg}
\IEEEauthorblockA{DD2430, KTH Royal Institute of Technology\\
Project in collaboration with Ericsson}
}

\maketitle

\begin{abstract}
% 150-200 words, written last: problem, approach, main quantitative result, main limitation.
\todo{Write last. One sentence each: why model lowering matters, what we built,
how we evaluated it, the headline result with its noise band, the main limitation.}
\end{abstract}

\begin{IEEEkeywords}
compilers, XLA, JAX, large language models, agents, performance optimization
\end{IEEEkeywords}

% ---------------------------------------------------------------------------
\section{Introduction}
% About 1 page.
% - Running AI models is costly, so speeding up inference saves money.
% - Lowering pipeline: experts redesign algorithms, engineers write kernels,
%   compilers fuse and schedule automatically. Compilers use fixed rewrite rules.
% - Idea: an LLM agent that reads compiler output and profiles, rewrites code,
%   and measures the result, like a human performance engineer.
% - Scope: inference (forward pass) on one NVIDIA L4 GPU.
% - Contributions as a short bulleted list.
\todo{Motivation and problem in two paragraphs.}

Our contributions are:
\begin{itemize}
  \item \todo{The agent loop around JAX/XLA: observe, propose, verify, measure.}
  \item \todo{An evaluator that checks correctness before timing, and a noise-floor measurement.}
  \item \todo{An evaluation on MLP forward passes on an NVIDIA L4, against vanilla XLA.}
\end{itemize}

% ---------------------------------------------------------------------------
\section{Related Work}
% About 0.75 page. Compare each line of work to ours.
\subsection{Hand-designed algorithmic rewrites}
\todo{FlashAttention~\cite{dao2022flashattention} as an example of a rewrite a compiler cannot find.}

\subsection{Machine-learning compilers}
\todo{JAX~\cite{jax} and XLA~\cite{openxla}: fusion, layout, scheduling, and what they cannot change.}

\subsection{LLM agents for kernel and code optimization}
\todo{KernelBench~\cite{ouyang2025kernelbench}, the Sakana AI CUDA Engineer
work~\cite{lange2025robust}, AlphaEvolve \todo{add citation}. How ours differs:
JAX-level rewrites, compiler IR and profiles as feedback.}

% ---------------------------------------------------------------------------
\section{Problem Description}
% About 0.5 page. Make the task precise enough that the results can be judged.
% - Input: a pure JAX function f and example inputs x.
% - Output: a function g with g(x) close to f(x) within tolerance, and lower latency.
% - Correctness: |g(x) - f(x)| <= atol + rtol * |f(x)|, with the values used.
% - Objective: median latency on the target device, excluding compilation.
% - Baseline: vanilla XLA compilation of f.
\todo{Formal statement of the optimization task.}

\begin{equation}
\text{speedup} = \frac{t_{\text{median}}(f)}{t_{\text{median}}(g)}
\quad \text{subject to} \quad
|g(x) - f(x)| \le \text{atol} + \text{rtol}\,|f(x)|
\end{equation}

\todo{State what counts as a real improvement: a speedup outside the noise band.}

% ---------------------------------------------------------------------------
\section{Methodology and Implementation}
% About 2 pages. One subsection per component, plus a system diagram.
\subsection{System overview}
\todo{The loop: observe (HLO + profile), propose (LLM), verify (evaluator), measure, repeat.}

\begin{figure}[t]
  \centering
  \fbox{\parbox{0.9\columnwidth}{\centering\vspace{2em}\todo{system diagram}\vspace{2em}}}
  \caption{\todo{The agent loop around the XLA compiler.}}
  \label{fig:system}
\end{figure}

\subsection{Compiler interface}
\todo{\texttt{compile\_program}: lowering to StableHLO, compiling to optimized HLO.}

\subsection{Profiler}
\todo{What is measured, how, and what the agent sees.}

\subsection{Evaluator}
\todo{Correctness before timing, tolerances, failure stages, median latency.}

\subsection{LLM policy and prompting}
\todo{Gemini model and version, prompt contents (HLO, profile, source, history),
JSON reply format, how candidate code is loaded, retries and timeouts.}

\subsection{Numerical precision on the GPU}
\todo{Default L4 matmul precision is low (TF32-like): median relative error
$2.9\times10^{-4}$ vs $2.4\times10^{-7}$ with \texttt{highest}. All runs use \texttt{highest}.}

% ---------------------------------------------------------------------------
\section{Experimental Evaluation}
% About 2.5 pages. Setup first, then each result with its noise band.
\subsection{Setup}
\todo{Hardware (NVIDIA L4, g2-standard-4), software versions (JAX 0.12.0.1, CUDA 12),
matmul precision, Gemini model, repeats, seeds.}

\subsection{Benchmarks}
\todo{Seeded ReLU MLPs: 1--10 hidden layers, widths 32/128/256, 1000 examples in batches of 32.
Two reference programs: a scan-based version with a planted inefficiency, and a vectorized version
that is already well written.}

\subsection{Measurement noise}
\todo{Noise floor: identical programs evaluated against each other. Report the 5th--95th percentile band.}

\subsection{Agent results}
\begin{table}[t]
  \centering
  \caption{Agent speedup over vanilla XLA on the NVIDIA L4. \todo{fill in}}
  \label{tab:agent}
  \begin{tabular}{@{}lrrr@{}}
    \toprule
    Reference & Correct candidates & Median speedup & Above noise band \\
    \midrule
    Scan (planted) & \todo{} & \todo{} & \todo{} \\
    Vectorized & \todo{} & \todo{} & \todo{} \\
    \bottomrule
  \end{tabular}
\end{table}
\todo{Results per architecture (figure), example rewrite, failure cases.}

\subsection{JAX on GPU versus NumPy on CPU}
\todo{Optional context table. Note that it calls the program once per batch, so it
mostly measures launch overhead on the GPU.}

% ---------------------------------------------------------------------------
\section{Summary and Discussion}
% About 1 page.
% - What worked, what did not, and why.
% - Limitations: one model family, one GPU, one LLM, single-input correctness check,
%   no custom kernels, peak memory not measured.
% - Future work.
\todo{Summary of findings.}
\todo{Limitations and threats to validity.}
\todo{Future work.}

% ---------------------------------------------------------------------------
\section{Summary of Discussions with the Opposing Group}
% About 0.5 page. Fill in after the opposition on Dec 10.
\todo{Main points raised by the opposing group and how the report was changed in response.}

\section*{Acknowledgment}
We thank Alexander Kravberg at Ericsson for supervising the project.

\bibliographystyle{IEEEtran}
\bibliography{references}

\end{document}
